\chapter{Resultados y Discusión}

\section{Resultados de clasificación y de la segmentación complementaria}

\subsection{Resultados disponibles de clasificación}

\noindent\textit{Nota sobre consolidación.} Esta sección reemplaza una lectura v7 anterior (medias de tres semillas con remuestreo por imagen) por la lectura consolidada más reciente del banco completo (cómputo canónico, agrupado por \texttt{lesion\_id}), consistente con los números usados en el resto de este capítulo.

La tabla \ref{tab:v7-clasificacion} resume la evaluación congelada (el backbone no actualiza pesos; solo se ajusta la regresión logística) sobre la partición de HAM10000 (n = 2.004, agrupada por lesión). B71 es el brazo propuesto (SimSiam anclado con término de tono); A71 es la ablación sin ese término; el público usa el mismo protocolo de extracción que B71.

\begin{table}[htbp]
\centering
\caption{Clasificación congelada en HAM10000 (n = 2.004): accuracy, exactitud balanceada, F1 macro, AUC, ECE y MCC.}
\label{tab:v7-clasificacion}
\begingroup
\small
\setlength{\tabcolsep}{3pt}
\begin{tabular}{>{\raggedright\arraybackslash}p{4.1cm}rrrrrr}
\toprule
Brazo & Accuracy & Ex. balanceada & F1 macro & AUC OvR & ECE & MCC \\
\midrule
DINOv2 público (protocolo B71) & 0,8278 & 0,6659 & 0,6837 & 0,9619 & 0,0688 & 0,6712 \\
A71 (sin término de tono) & 0,8398 & 0,7305 & 0,7229 & 0,9644 & 0,0154 & 0,6996 \\
B71 (propuesto, con término de tono) & 0,8378 & 0,6964 & 0,7141 & 0,9602 & 0,0529 & 0,6901 \\
\bottomrule
\end{tabular}
\endgroup
\end{table}
\noindent\textit{Nota.} Elaboración propia.

Los intervalos de confianza individuales de A71 y B71 se solapan en exactitud balanceada y F1 macro. Esta comparación de intervalos individuales no demuestra una ventaja general del término de tono de B* sobre A*; una afirmación más fuerte requiere examinar el contraste pareado formal entre ambos brazos.

\subsection{Avance y límites del módulo complementario de segmentación}
\label{sec:segmentacion-complementaria}

La línea de segmentación sin etiquetas sigue un flujo separado de MoCo v2/ResNet-50, CCAM, pseudo-máscaras y un U-Net estudiante. El candidato S4-A2 obtuvo Dice 0,7221 en ISIC2018 y 0,7949 en HAM10000; superó por márgenes pequeños algunos umbrales registrados frente al estudiante \texttt{student\_r3}. Sin embargo, las pruebas de integración C1 no muestran una mejora general del clasificador al usar máscaras. Este resultado respalda mantener la segmentación como módulo complementario, no afirmar que mejore el clasificador.

\section{Discusión confirmatoria, ablaciones, adaptación de dominio y limitaciones}

\subsection{Comparación confirmatoria: clasificador SSL congelado frente a ocho líneas base}

\enquote{Congelado} significa que el backbone preentrenado con SSL no actualiza sus pesos: solo se ajusta una regresión logística sobre las representaciones extraídas, con etiquetas del conjunto de entrenamiento de cada dataset. Esta es la comparación confirmatoria y controlada del banco: todas las líneas base se evalúan sobre la misma partición de HAM10000 (n = 2.004, agrupada por \texttt{lesion\_id}), con bootstrap pareado de 1.000 remuestras, semilla 20260923.

\begin{table}[htbp]
\centering
\caption{Comparación confirmatoria: F1 macro del clasificador congelado (B71) contra ocho líneas base, HAM10000 (n = 2.004).}
\label{tab:confirmatoria-baselines}
\begin{tabular}{lrrl}
\toprule
Línea base & F1 macro propio & $\Delta$ F1 vs. B71 [IC 95\,\%] & Lectura \\
\midrule
VGG19 & 0,3810 & $+$0,3331 [$+$0,2844; $+$0,3800] & excluye cero \\
AlexNet & 0,4728 & $+$0,2413 [$+$0,1942; $+$0,2898] & excluye cero \\
ResNet-50 & 0,4670 & $+$0,2471 [$+$0,1969; $+$0,2943] & excluye cero \\
EfficientNet-B0 & 0,5041 & $+$0,2100 [$+$0,1579; $+$0,2609] & excluye cero \\
SimSiam ResNet-50 & 0,5369 & $+$0,1772 [$+$0,1278; $+$0,2299] & excluye cero \\
DINOv2 ViT-S/14 & 0,5508 & $+$0,1633 [$+$0,1198; $+$0,2066] & excluye cero \\
ViT-B/16 & 0,5633 & $+$0,1509 [$+$0,1005; $+$0,2068] & excluye cero \\
DINOv2 ViT-B/14 & 0,6002 & $+$0,1139 [$+$0,0728; $+$0,1585] & excluye cero \\
\bottomrule
\end{tabular}
\end{table}
\noindent\textit{Nota.} Elaboración propia.

B71 alcanza F1 macro 0,7141 en esta misma partición. Los ocho contrastes pareados excluyen cero: es el único resultado del banco con lectura confirmatoria 8 de 8.

\subsection{Comparación descriptiva: ajuste fino de las ocho líneas base}

El ajuste fino actualiza todos los pesos de cada línea base con etiquetas, usando una partición distinta (n = 2.014) y eligiendo la época por desempeño en validación, lo que introduce un sesgo optimista declarado. Por estas diferencias de protocolo, esta comparación es descriptiva y no un contraste pareado contra B71.

Los máximos de ajuste fino fueron 0,7058 (ViT-B/16) en HAM10000, 0,7962 (ViT-B/16) en ISIC2019 y 0,5889 (VGG19) en ISIC2020. A diferencia del régimen congelado, donde B71 gana a las ocho líneas base de forma confirmatoria, el ajuste fino de los mejores backbones sí alcanza o supera a B71 en HAM10000 e ISIC2019. La distinción entre ambos regímenes --representación congelada frente a ajuste con etiquetas-- es la que sostiene la contribución SSL: B71 no requiere etiquetas para aprender la representación, mientras que el ajuste fino las requiere en su totalidad.

\subsection{Caso especial: ISIC2020, F1 contra sensibilidad clínica}

En ISIC2020, el contraste pareado entre B71-limpio y sus ocho ajustes finos, corregido por Holm, favorece a B71 en F1 macro frente a 7 de 8 líneas base (VGG19 queda cerca, $p_{\text{Holm}}=0{,}066$). Sin embargo, la exactitud balanceada favorece a 7 de 8 ajustes finos, y la sensibilidad de melanoma --la métrica de mayor relevancia clínica-- muestra la brecha más marcada del banco: B71 detecta 39 de 122 melanomas (recall 0,3197), frente a 86 de 122 (recall 0,7049) del AlexNet ajustado. F1 macro alto no implica mejor detección de melanoma; esta tensión se reporta de forma explícita y no se resuelve a favor de ninguno de los dos regímenes sin una decisión clínica adicional sobre qué métrica priorizar.

\subsection{Ablación de segmentación complementaria (F4)}

F4 concatena la imagen completa, un recorte de la lesión con margen del 10\,\% y descriptores de forma y color de la pseudo-máscara del segmentador (\texttt{student\_r3}); no aplica enmascaramiento duro. En el contraste canónico B71+F4 frente a B71 en HAM10000, $\Delta$F1 macro $=+0{,}0300$ [$-0{,}0046$; $+0{,}0715$], exactitud balanceada $+0{,}0222$ [$-0{,}0159$; $+0{,}0624$] y recall de melanoma $+0{,}0135$ [$-0{,}0446$; $+0{,}0714$]: ningún intervalo excluye cero. La sensibilidad por semilla confirma la falta de robustez: $+0{,}0300$ (semilla 42), $+0{,}0092$ (semilla 43) y $+0{,}0261$ (semilla 44), con 0 de 3 intervalos excluyendo cero. En el análisis exploratorio E3b, F4 sube el punto de F1 en 8 de 8 backbones, pero solo 5 de 8 intervalos excluyen cero. F4 se mantiene como variante exploratoria; B71 sin segmentación sigue siendo la referencia principal del banco.

\subsection{Adaptación de dominio: continuación SSL con PAD-UFES-20}

D2 y D3 parten del mismo backbone limpio (A* semilla 42) y continúan preentrenamiento SSL por 2.000 pasos adicionales: D2 con el pool fuente sin cambios, como control de cómputo; D3 agrega 1.141 imágenes de la partición \texttt{target\_adapt} de PAD-UFES-20 sin etiquetas (aprox. 2,7\,\% del pool total de D3). Esta continuación adapta las representaciones al dominio de fotografías clínicas. Sobre ambos se ajusta la misma regresión logística con etiquetas de HAM10000 (n = 8.011), congelada, y se evalúa en la partición \texttt{target\_test} de PAD-UFES-20 (solo BCC, melanoma y nevo; 567 de 1.142 imágenes, 363 pacientes, 25 melanomas). Las imágenes de \texttt{target\_test} no participan en la adaptación SSL ni en el ajuste supervisado; sus etiquetas se utilizan para calcular las métricas de evaluación.

\begin{table}[htbp]
\centering
\caption{Adaptación de dominio a PAD-UFES-20: comparación D0--D3.}
\label{tab:adaptacion-pad}
\begin{tabular}{lrrrr}
\toprule
Brazo & F1 macro & Exactitud balanceada & MCC & Fuera de rango \\
\midrule
D0 (DINOv2 público) & 0,4410 & 0,5284 & 0,3740 & 0,3210 \\
D1 (A* s42 limpio) & 0,2443 & 0,4244 & 0,1871 & 0,4039 \\
D2 (continuación, fuente) & 0,2877 & 0,4897 & 0,2309 & 0,3422 \\
D3 (D2 + PAD sin etiquetas) & 0,3693 & 0,5271 & 0,3039 & 0,4780 \\
\bottomrule
\end{tabular}
\end{table}
\noindent\textit{Nota.} Elaboración propia.

``Fuera de rango'' es la proporción de predicciones \texttt{argmax} fuera de BCC/MEL/NV.

El contraste primario D3$-$D2 es significativo: $\Delta$F1 $=+0{,}0816$ [$+0{,}0525$; $+0{,}1085$] y $\Delta$MCC $=+0{,}0730$ [$+0{,}0419$; $+0{,}1034$]. La exactitud balanceada ($+0{,}0375$ [$-0{,}0161$; $+0{,}0875$]) y el AUC ($+0{,}0162$ [$-0{,}0065$; $+0{,}0393$]) no excluyen cero. Sin embargo, D3 sigue por debajo del encoder público D0: pierde $-0{,}0717$ [$-0{,}1100$; $-0{,}0332$] en F1 y $-0{,}0701$ [$-0{,}1188$; $-0{,}0200$] en MCC. Los recalls por clase de D3 son BCC 0,302, melanoma 0,320 y NV 0,959, y el 47,8\,\% de las predicciones caen fuera de BCC/MEL/NV. La lectura es una señal piloto frente al control de cómputo D2, no evidencia de que la adaptación a PAD esté resuelta: una sola semilla, muy pocos melanomas de prueba y un mapeo parcial de diagnósticos.

\subsection{Limitaciones de ITA y Fitzpatrick}

El estimador de ITA del proyecto no es un clasificador Fitzpatrick validado: acierta 22,9\,\% frente al 51,3\,\% de una base de referencia por mayoría en PAD-UFES-20 \parencite{alipour2026italimitations}. Los metadatos de PAD-UFES-20 no permiten concluir sobre fototipos V--VI: esos grupos tienen muy pocas observaciones en el banco. Por estas razones, ningún resultado de este capítulo se presenta como una auditoría concluyente de equidad por fototipo.

\subsection{Estado del objetivo y trabajo pendiente}

El banco respalda el aprendizaje autosupervisado como método de representación frente a ocho líneas base, de forma confirmatoria en régimen congelado (8 de 8), y aporta una señal piloto de que la adaptación con imágenes de PAD sin etiquetas mejora la sonda de origen frente a su control de cómputo. No demuestra que la segmentación (F4) mejore la clasificación de forma general, que el ITA prediga Fitzpatrick, ni que la adaptación a PAD supere al encoder público. Quedan pendientes: (i) fijar antes de nuevas semillas el split y los hashes de D2/D3, hoy sin \texttt{data\_revision} registrada; (ii) separar el recorte de los descriptores en F4 para aislar el efecto de la segmentación; (iii) validar el ITA contra anotación experta antes de cualquier afirmación de equidad; y (iv) decidir, con criterio clínico, si la detección de melanoma (recall) o el F1 macro agregado es la métrica que debe gobernar la selección final del modelo, dada la tensión observada en ISIC2020. Ninguna de estas tareas autoriza por sí sola nuevas corridas; cualquier corrida se ejecuta en UNAB con presupuesto previamente acordado.
